\chapter*{全书结构}

本卷按 Compute Systems Engine 的工程演化组织，但不再把一个小动作拆成一个目录项。正文每章控制在三到四个大节：先给具体问题和输入样例，再推导类型与接口，接着落到实现、测试和探针，最后给阶段验收。细节仍然要讲透，只是放在大节内部展开。

本卷现在有两条可运行 lab 主线。第一条是 \filepath{reference_pipeline}，从 \code{event,value} 输入开始，建立输入身份、reference、结构化诊断、report、manifest、split、partial result 和 merge。第二条是 \filepath{compute_systems}，承接第二册原理卷，为 CPU、缓存、TLB、NUMA、SIMD、线程调度、同步、原子、任务运行时、异步 I/O、分布式 quorum 和 checkpoint/recovery 提供逐章模型、探针和测试。

每个阶段都必须产出四类东西。第一，稳定接口，例如 \code{InputSplit}、\code{RecordId}、\code{HotBatch}、\code{QueueContract}、\code{ManifestEntry} 或 \code{RunReport}。第二，可运行实现，至少包含 reference 和一个带机制的版本。第三，测试，包括正常输入、边界输入和能击穿朴素方案的失败输入。第四，报告，记录正确性、性能、资源、故障和未验证边界。

本卷与第二册原理卷并行阅读。原理卷先解释机制为什么存在，本卷再把机制落成代码。若读者只读本卷而跳过原理卷，容易把接口当成固定答案；若只读原理卷而不做本卷，容易把工程边界停留在口头推导。两卷合起来，才形成从零到精通的完整训练。
